\section{Introduction}
\label{sec:introduction}

Ghana's external sector is highly exposed to commodity markets. UN Trade and Development classifies an economy as commodity-dependent when commodities account for more than 60\% of merchandise exports, and its 2025 country profile reports that Ghana's three leading commodity exports represented 78.4\% of allocated product exports over 2021--2023. Bank of Ghana data also underline the continuing importance of cocoa, gold and crude oil: in 2025, exports of cocoa beans and products were about US\$4.0 billion, non-monetary gold about US\$21.0 billion, and crude oil about US\$2.6 billion \citep{unctad2025,bog2025}. These exposures make the interaction between international commodity prices and the Ghanaian cedi relevant for financial risk measurement, portfolio decisions and the assessment of extreme market conditions.

The Ghana-specific literature already shows that commodity--currency relationships are nonlinear and time-varying. \citet{archer2022} report asymmetric relationships between cocoa, gold, Brent crude oil and the exchange rate across quantiles. \citet{barson2023} find that commodity-return comovement varies across horizons and that the exchange rate plays an important role in those relationships. \citet{yeboah2025} similarly document nonlinear and asymmetric relationships between commodity prices and Ghanaian macroeconomic variables. Daily-frequency work has also linked the cedi to global uncertainty and broader cross-market connectedness \citep{atipaga2025,woode2025}. Copula methods are not new in Ghana: \citet{mensah2020} estimate static and time-varying copulas for major foreign currencies against the cedi and report tail-dependent exchange-rate comovement.

International evidence likewise suggests that commodity--currency dependence varies with market state, commodity composition and model choice. Oil and exchange-rate comovement can strengthen after crises \citep{reboredo2012}, while time-varying tail dependence has been reported for major commodity currencies \citep{ignatieva2017}. Recent studies document regime dependence, asymmetric risk transmission and crisis-related changes in commodity and currency markets \citep{qiao2023,go2024,ning2025,dodd2026}. These results motivate direct analysis of extreme dependence, but they do not imply that every pair or every crisis window must exhibit stronger tail dependence.

Two methodological issues are especially important for the present setting. First, commodity futures and the cedi do not necessarily share the same observation dates. Classical nonsynchronous-trading research shows that mismatched observation times can distort return moments and cross-market dependence \citep{scholes1977,lo1990,martens2001,schotman2006}. Calendar construction is therefore part of the statistical specification rather than a neutral data-cleaning step. Second, relative model selection does not establish absolute adequacy. A lowest-AIC copula can still be misspecified, and misspecified marginal distributions can propagate error into copula and downstream risk estimates \citep{fantazzini2009,genest2008,genest2009}.

This study addresses those issues using two co-equal constructions of the same daily-source data for cocoa, gold, Brent crude oil, West Texas Intermediate (WTI) crude oil and the Ghanaian cedi over 2015--2026. Panel A imposes a business-day calendar with forward filling capped at three business days before returns are computed. Panel B uses only consecutive dates on which all five source series have actual observations and therefore contains no forward-filled return observations. Neither panel is designated as the primary or corrected dataset. A conclusion is described as calendar-robust when it survives both constructions and calendar-sensitive when its inferential classification changes.

The paper makes three contributions. First, it evaluates eight static bivariate copula families with an absolute Rosenblatt-transform Cram\'{e}r--von Mises goodness-of-fit procedure and parametric bootstrap rather than equating relative fit with correctness. Second, it compares empirical lower- and upper-tail concentration across calm periods and prespecified COVID-19 and 2024 stress windows using a moving-block bootstrap, with Bonferroni and Benjamini--Hochberg adjustment across each 60-test family. Third, it treats the cedi's zero-heavy daily return distribution and marginal-model inadequacy explicitly, so commodity--cedi copula and extreme-value quantities are reported as conditional diagnostics rather than model-free structural parameters.

The resulting research question is deliberately narrow: which conclusions about commodity tail dependence and exchange-rate risk remain defensible after calendar construction, marginal adequacy, absolute copula fit and multiple testing are made explicit? This emphasis on robustness is important for applied international finance because decisions based on tail dependence can be sensitive not only to the chosen model, but also to the observation process that generates the return series.
